\section{Quotient metrics, local indices, and the divisor inequality}
\label{div:section}

We now prove the determinant and lattice assertions underlying the divisor
inequality. In particular, the inequality between a quotient covolume and
finite Gram determinants is derived, rather than taken as a hypothesis.
We also prove the operator estimates for a chain of rectangular matrices.
The source assertions being reconstructed occur in the quotient-metric
and local-to-global divisor constructions.\cite{alpo-uniform-source}
These results do not yet construct the particular automorphic lattices or
establish their uniform estimates. The remaining inputs are described at
the end of the section.

\subsection{Normalization and local indices}

Let $K$ be a totally imaginary number field, let $d=[K:\mathbb Q]$,
and put $r_2=d/2$. Write $\mathcal O=\mathcal O_K$ and let $v$ run over
the complex places, one from each pair of conjugate embeddings. For a
nonzero prime ideal $\mathfrak p$ put
\[
 w_{\mathfrak p}=\frac{\log N\mathfrak p}{d},\qquad
 \operatorname{ord}_{\mathfrak p}:K^\times\longrightarrow\mathbb Z.
\]
Valuations in this section are applied only to nonzero elements; a zero
determinant is never assigned a finite valuation. The product formula is
\begin{equation}\label{div:product-formula}
 \sum_{\mathfrak p}w_{\mathfrak p}\operatorname{ord}_{\mathfrak p}(a)
 =\frac2d\sum_{v\mid\infty}\log|v(a)|\qquad(a\in K^\times).
\end{equation}
Indeed, both sides multiplied by $d$ are
$\log|N_{K/\mathbb Q}(a)|$, computed from the factorization of the
principal fractional ideal and from the complex embeddings, respectively.

A full lattice in $K^q$ means a finitely generated $\mathcal O$-submodule
$L$ spanning $K^q$. Its determinant ideal $\mathfrak a(L)$ is the
fractional ideal generated by
$\det(x_1,\ldots,x_q)$ with $x_i\in L$, in the displayed coordinates.
Put $a_{\mathfrak p}(L)=\operatorname{ord}_{\mathfrak p}\mathfrak a(L)$.
For positive definite Hermitian matrices $G_v$ on $\mathbb C^q$, define
\begin{equation}\label{div:covolume}
 C_G(L)=N\mathfrak a(L)\prod_{v\mid\infty}\det G_v,
 \qquad c_G(L)=\frac2d\log C_G(L).
\end{equation}
Here the norm of a fractional ideal is positive and multiplicative.
The first expression is the real Minkowski covolume divided by
$2^{-r_2q}|\Delta_K|^{q/2}$, when the metric at a complex place is
the real part of its Hermitian form. Thus it removes exactly the
covolume of $\mathcal O^q$ with its standard metrics. The formula follows
first for a free lattice from the real determinant of a complex matrix,
which is the squared absolute value of its complex determinant, and then
for an arbitrary lattice by passing to a finite-index free sublattice.
Changing coordinates transforms the determinant ideal and the Gram
matrices by reciprocal factors; \eqref{div:product-formula} shows that
$C_G(L)$ is unchanged.

The usual normalized arithmetic degree and slope in this convention are
\[
 \widehat{\deg}(L)=-\tfrac12c_G(L),\qquad
 \mu(L)=-\frac{c_G(L)}{2q}\quad(q>0).
\]
This specifies the factor of two: a convention using logarithms of
squared norms has slope $2\mu(L)$. For the ordinary linear dual lattice
$L^\vee=\operatorname{Hom}_{\mathcal O}(L,\mathcal O)$ with dual metrics,
the determinant ideal and metric determinants are inverted. Consequently
\begin{equation}\label{div:ordinary-dual}
 c(L^\vee)=-c(L),\qquad 2\mu(L^\vee)=c(L)/q.
\end{equation}

For full lattices $A,B$ in the same $K$-space define the signed local index
\[
 i_{\mathfrak p}(A,B)=
 \operatorname{length}_{\mathcal O_{\mathfrak p}}
       \bigl(A_{\mathfrak p}/(A_{\mathfrak p}\cap B_{\mathfrak p})\bigr)
 -\operatorname{length}_{\mathcal O_{\mathfrak p}}
       \bigl(B_{\mathfrak p}/(A_{\mathfrak p}\cap B_{\mathfrak p})\bigr).
\]
The subscript denotes ordinary localization at the single prime.

\begin{lemma}[Index and covolume]\label{div:index}
For full lattices $A,B$ and fixed metrics $G$,
\begin{align}
 i_{\mathfrak p}(A,B)&=a_{\mathfrak p}(B)-a_{\mathfrak p}(A),
       \label{div:index-det}\\
 c_G(B)-c_G(A)&=2\sum_{\mathfrak p}w_{\mathfrak p}
                                  i_{\mathfrak p}(A,B).
       \label{div:index-covolume}
\end{align}
All sums are finite. The indices are additive:
$i_{\mathfrak p}(A,C)=i_{\mathfrak p}(A,B)+i_{\mathfrak p}(B,C)$.
\end{lemma}
\begin{proof}
Over the discrete valuation ring $\mathcal O_{\mathfrak p}$, choose a
basis of $A_{\mathfrak p}$. Elementary divisors give a basis in which
$B_{\mathfrak p}=\bigoplus_{j=1}^q\pi^{b_j}\mathcal O_{\mathfrak p}e_j$,
with $b_j\in\mathbb Z$. The two lengths in the definition are
$\sum_j\max(b_j,0)$ and $\sum_j\max(-b_j,0)$. Their difference is
$\sum_jb_j$, the valuation of the change-of-basis determinant. This
proves \eqref{div:index-det}. The determinant ideals have finite support,
and taking logarithms in \eqref{div:covolume} proves
\eqref{div:index-covolume}. Additivity follows from
\eqref{div:index-det}.
\end{proof}

\begin{lemma}[Determinantal ideals of rectangular matrices]\label{div:minors}
Let $A\in\operatorname{Mat}_{q\times n}(K)$ have rank $q$. Then
\[
 a_{\mathfrak p}(A\mathcal O^n)
 =\min_{\substack{I\subset\{1,\ldots,n\}\\|I|=q}}
           \operatorname{ord}_{\mathfrak p}\det A_I,
\]
where zero minors are omitted. More generally, for compatible matrices
$A,B$ and every $k$, the ideal generated by the $k$-minors of $AB$ is
contained in the product of the ideals generated by the $k$-minors of
$A$ and of $B$.
\end{lemma}
\begin{proof}
The exterior products of $q$ columns generate the top exterior power of
the image module, so its determinant ideal is precisely the ideal
generated by the maximal minors. An ideal over a discrete valuation
ring has valuation equal to the minimum valuation of any finite set of
nonzero generators. For the second assertion, Cauchy--Binet expresses
each $k$-minor of $AB$ as a sum of products of a $k$-minor of $A$ and a
$k$-minor of $B$. This proves the ideal inclusion, including cases in
which some minors vanish.
\end{proof}

\subsection{Hermitian duals and simultaneous localization}

Suppose that $K$ has an involution $a\mapsto\bar a$. This bar need not
act as complex conjugation at every embedding. Put
$\overline{\mathfrak p}=\{\bar a:a\in\mathfrak p\}$; then
$N\overline{\mathfrak p}=N\mathfrak p$ and
$\operatorname{ord}_{\mathfrak p}\bar a
 =\operatorname{ord}_{\overline{\mathfrak p}}a$.
Use $A^*=\bar A^{,t}$ for matrices over $K$ and ordinary conjugate
transpose for matrices over $\mathbb C$.

The ring relevant to a Hermitian local basis is
\[
 \mathcal O_{\{\mathfrak p,\overline{\mathfrak p}\}}
 =\mathcal O_{\mathfrak p}\cap
       \mathcal O_{\overline{\mathfrak p}}\subset K.
\]
It is obtained by inverting elements outside both primes and is stable
under the involution. A simultaneous local basis of $L$ at
$\mathfrak p,\overline{\mathfrak p}$ is a $K$-basis $e$ which is a basis
of both $L_{\mathfrak p}$ and $L_{\overline{\mathfrak p}}$.
Such a basis always exists. To see this, choose bases of the reductions
of the projective module $L$ modulo the two primes, choosing the same
basis if the primes coincide. The Chinese remainder theorem lifts these
to $q$ vectors of $L$. At each selected prime their determinant is a unit
in a local basis, so they form a basis there. This also proves their
$K$-linear independence. Equivalently, every finite projective module
of constant rank over this semilocal ring is free.

Let $H=H^*\in\operatorname{GL}_q(K)$, write
$h_H(x,y)=\bar x^{,t}Hy$, and define the Hermitian dual
\[
 L^\#=\{y\in K^q:h_H(x,y)\in\mathcal O\text{ for every }x\in L\},
 \qquad j_{\mathfrak p}(H,L)=i_{\mathfrak p}(L,L^\#).
\]
For a basis matrix $E$ with columns $e_1,\ldots,e_q$, its Gram
determinant is $g_H(e)=\det(E^*HE)\in K^\times$.

\begin{lemma}[The Hermitian index]\label{div:hermitian-index}
If $e$ is a simultaneous local basis of $L$ at
$\mathfrak p,\overline{\mathfrak p}$, then
\begin{align}
 \operatorname{ord}_{\mathfrak p}g_H(e)
 &=\operatorname{ord}_{\mathfrak p}\det H
       +a_{\mathfrak p}(L)+a_{\overline{\mathfrak p}}(L),
       \label{div:gram-valuation}\\
 j_{\mathfrak p}(H,L)&=-\operatorname{ord}_{\mathfrak p}g_H(e).
       \label{div:j-gram}
\end{align}
In particular $j_{\overline{\mathfrak p}}(H,L)=j_{\mathfrak p}(H,L)$.
\end{lemma}
\begin{proof}
The first identity is the determinant identity
$g_H(e)=\overline{\det E}\det H\det E$, using that $E$ represents $L$
at both primes. If $L^{\mathrm{lin},\vee}$ denotes the lattice dual
under the ordinary bilinear dot product, then
$L^\#=H^{-1}(\bar L)^{\mathrm{lin},\vee}$. It follows that
\[
 a_{\mathfrak p}(L^\#)
 =-\operatorname{ord}_{\mathfrak p}\det H
       -a_{\overline{\mathfrak p}}(L).
\]
Apply \eqref{div:index-det}. Finally $\det H$ is fixed by the
involution, so the right side of \eqref{div:gram-valuation} is invariant
under $\mathfrak p\leftrightarrow\overline{\mathfrak p}$.
\end{proof}

If $v(\bar H)=\overline{v(H)}$ for every embedding and every $v(H)$ is
positive definite, then the metrics $G_v=v(H)$ give the exact identity
\begin{equation}\label{div:rational-hermitian}
 c_H(L)=\sum_{\mathfrak p}w_{\mathfrak p}
                    \operatorname{ord}_{\mathfrak p}g_H(e_{\mathfrak p})
       =-\sum_{\mathfrak p}w_{\mathfrak p}j_{\mathfrak p}(H,L).
\end{equation}
Indeed, sum \eqref{div:gram-valuation}, use invariance of the weights
under conjugation, and apply \eqref{div:product-formula} to $\det H$.
The simultaneous bases may vary with $\mathfrak p$; the resulting
valuations do not.

\subsection{The dual exponent in the finite-group construction}

We can also derive the integer controlling the ambient dual lattice.
This uses the actual finite class sets, stabilizers and character
projectors occurring in the construction, rather than a bound for an
unspecified lattice.

\begin{lemma}[Dual of a saturated section]\label{div:cut-dual}
Let $(V,h)$ be a nondegenerate Hermitian $K$-space, let $N\subset V$
be a full lattice, and let $W\subset V$ be a nondegenerate subspace.
Let $\pi:V\to W$ be the orthogonal projection and put $N_W=N\cap W$.
For the restricted form on $W$ one has
\[
                       N_W^\#=\pi(N^\#).
\]
For the form $h/a$ on $W$, with $a\in\mathbb Q_{>0}$, the dual is
$a\pi(N^\#)$.
\end{lemma}
\begin{proof}
The inclusion from right to left follows from
$h(x,\pi z)=h(x,z)$ for $x\in W$.
For the reverse inclusion, $N_W$ is saturated in $N$: if a nonzero
$b\in\mathcal O$ satisfies $bx\in W$, then $x\in W$.
The quotient $N/N_W$ is consequently a finitely generated torsion-free
module over the Dedekind domain $\mathcal O$, hence projective. The
exact sequence splits. Given $y\in N_W^\#$, the conjugate-linear
functional $x\mapsto h(x,y)$ on $N_W$ extends, by zero on a chosen
complement, to an $\mathcal O$-valued conjugate-linear functional on
$N$. Extending scalars to $K$ and using nondegeneracy of $h$ gives a
vector $z\in N^\#$ which represents it. Thus
$h(x,\pi z-y)=0$ for all $x\in N_W$, and hence for all $x\in W$,
because $N_W$ spans $W$. Nondegeneracy on $W$ gives $\pi z=y$.
Finally $h(x,y)/a=h(x,y/a)$, proving the scaling assertion.
\end{proof}

For this construction fix an embedding $v_0:K\hookrightarrow\mathbb C$
which intertwines the involution with complex conjugation.
Let $\mathcal C$ be a finite set. For each $c\in\mathcal C$, let
$\Gamma_c$ be a finite group acting by isometries on a Hermitian
$K$-space $(V_c,h_c)$ which is positive definite at the distinguished
embedding $v_0$. Let $N_c,D_c$ be full lattices such that
\begin{equation}\label{div:class-lattices}
 D_c\subset N_c,\qquad \Gamma_cD_c=D_c,\qquad
                         sN_c^\#\subset D_c
\end{equation}
for one positive integer $s$. Form the space and its class-sum form
\[
 \mathcal V=\bigoplus_{c\in\mathcal C}V_c^{\Gamma_c},\qquad
 h(x,y)=\sum_{c\in\mathcal C}\frac{h_c(x_c,y_c)}{|\Gamma_c|},
\]
and the full lattices
\[
 N=\bigoplus_c(N_c\cap V_c^{\Gamma_c}),\qquad
 D'=\bigoplus_c(D_c\cap V_c^{\Gamma_c}).
\]
The restriction to each fixed subspace is nondegenerate, because it
is positive definite at $v_0$.

\begin{proposition}[Class-sum dual containment]\label{div:class-dual}
For the preceding form and lattices, $sN^\#\subset D'\subset N$.
\end{proposition}
\begin{proof}
The operator $\pi_c=|\Gamma_c|^{-1}\sum_{\gamma\in\Gamma_c}\gamma$
is the orthogonal projection onto $V_c^{\Gamma_c}$: it is idempotent,
its image is the fixed subspace, and its adjoint equals itself because
$\gamma^*=\gamma^{-1}$. Lemma~\ref{div:cut-dual} gives
\[
 (N_c\cap V_c^{\Gamma_c})^\#_{h_c/|\Gamma_c|}
          =|\Gamma_c|\pi_c(N_c^\#).
\]
Multiplying by $s$, the right side becomes
$\sum_{\gamma\in\Gamma_c}\gamma(sN_c^\#)$, which lies in $D_c$
by \eqref{div:class-lattices} and is fixed by $\Gamma_c$.
Taking the orthogonal direct sum proves the assertion. Notice that
no stabilizer-order factor survives in the result: the class-sum
weight cancels the denominator in the averaging projection.
\end{proof}

Suppose a finite group $G$ acts on $\mathcal V$ by isometries and
preserves $D'$. Let $\mathcal X$ be a finite set of distinct irreducible
complex characters of $G$ whose values belong to $K$, identified with
its image under $v_0$, let $P\subset G$ be a
subgroup, and let $\psi:P\to K^\times$ be a character. In $K[G]$ put
\begin{align*}
 e_{\mathcal X}
   &=\sum_{\chi\in\mathcal X}\frac{\chi(1)}{|G|}
                      \sum_{g\in G}\chi(g^{-1})g,\\
 e_{P,\psi}&=\frac1{|P|}\sum_{g\in P}\psi(g)^{-1}g,
       &e&=e_{\mathcal X}e_{P,\psi}.
\end{align*}
These formulas act on $\mathcal V$ through the specified representation.

\begin{lemma}[Integral denominators of the projectors]\label{div:projectors}
The operator $e$ is a self-adjoint idempotent. The coefficients of
$|G||P|e$ lie in $\mathcal O$. Its image is the intersection of the
sum of the $\chi$-isotypic subspaces, $\chi\in\mathcal X$, with the
subspace on which $P$ acts through $\psi$.
\end{lemma}
\begin{proof}
After extending scalars through $v_0$ to $\mathbb C$, for an irreducible
representation with character $\vartheta$,
the central element
$e_\chi=\chi(1)|G|^{-1}\sum_g\chi(g^{-1})g$ acts by a scalar,
by Schur's lemma. Taking its trace and using character orthogonality
shows that this scalar is $1$ for $\vartheta=\chi$ and $0$ otherwise.
One can obtain the needed orthogonality directly: averaging
$A\mapsto\rho_\vartheta(g)A\rho_\chi(g)^{-1}$ on the space of
homomorphisms between the two representations
is the projection onto the intertwining maps. Its trace is
$|G|^{-1}\sum_g\vartheta(g)\chi(g^{-1})$, and its rank is one or
zero by Schur's lemma. Complete reducibility follows by averaging a
projection onto any invariant subspace. This proves that the $e_\chi$
are mutually orthogonal central idempotents.

The finite sum for $e_{P,\psi}$ directly shows both idempotence and
that its image is the $\psi$-eigenspace. It commutes with the central
operator $e_{\mathcal X}$. At $v_0$, finite-group representations can
be made unitary by averaging a positive form. Hence
$\overline{\chi(g)}=\chi(g^{-1})$, as an identity in $K$, and
$\overline{\psi(g)}=\psi(g)^{-1}$. Since the given representation
is already unitary for $h$, the two idempotents are self-adjoint.
Their commuting product is the orthogonal projection onto the
intersection of their images.

Every $\chi(g)$ is an algebraic integer: the eigenvalues of a
finite-order matrix are roots of unity. The same is true of
$\psi(g)^{\pm1}$, and $\chi(1)$ is an integer.
Thus $|G|e_{\mathcal X}$ and $|P|e_{P,\psi}$ have integral
coefficients. Multiplication proves the asserted denominator bound.
\end{proof}

\begin{theorem}[The actual ambient dual exponent]\label{div:ambient-exponent}
Put $W=e\mathcal V$, $M=N\cap W$, and equip $W$ with the scaled form
$h_W=h|_W/|P|$. Then
\begin{equation}\label{div:ambient-containment}
                         |G|sM^\#\subset M.
\end{equation}
For any further algebraic quotient of $(W,h_W,M)$ as in
Proposition~\ref{div:local-bound}, its local Gram determinants therefore
satisfy
\[
 \operatorname{ord}_{\mathfrak p}g_\Gamma(e_{\mathfrak p})
             \le D\operatorname{ord}_{\mathfrak p}(|G|s).
\]
Consequently, for any finite set $S$ of primes,
\[
 \sum_{\mathfrak p\in S}w_{\mathfrak p}
             \operatorname{ord}_{\mathfrak p}g_\Gamma(e_{\mathfrak p})/D
                         \le\log(|G|s)\qquad(D>0).
\]
\end{theorem}
\begin{proof}
Lemma~\ref{div:cut-dual} gives $M^\#=|P|e(N^\#)$ for the scaled
form. Proposition~\ref{div:class-dual} and
Lemma~\ref{div:projectors} imply
\[
 |G|sM^\#=|G||P|e(sN^\#)\subset |G||P|e(D')\subset D'\subset N.
\]
The left side also lies in $W$, proving
\eqref{div:ambient-containment}. Apply
Proposition~\ref{div:local-bound}. Finally $|G|s$ is a positive
integer, all its finite valuations are nonnegative, and
\eqref{div:product-formula} gives
$\sum_{\mathfrak p}w_{\mathfrak p}
\operatorname{ord}_{\mathfrak p}(|G|s)=\log(|G|s)$.
This bounds every partial sum.
\end{proof}

These objects also admit the rational Gram matrices required above.
If the entries of the initial Gram matrices and representation matrices
lie in a subfield $K_0$ on which every embedding respects conjugation,
first enlarge $K_0$ by the finitely many character values and roots of
unity used in the projectors. These elements have the same property:
after any embedding
they are still finite-group character values and roots of unity, so
conjugation is inversion on their eigenvalues. The enlarged subfield
still has this property, since it is preserved by field operations.
The equations for the fixed subspaces and the projector matrices now
have entries in $K_0$. Gaussian elimination
therefore supplies bases over $K_0$, and restriction of the forms in
those bases has entries in $K_0$. Positivity at every embedding is
preserved under restriction to an independent set of columns and under
division by the positive integer $|P|$. This proves the requisite
properties of the resulting ambient matrices once they hold for the
initial coefficient forms and actions.

\begin{example}[A coefficient normalization cannot be omitted]
Let $K=\mathbb Q(i)$ and let
$N=\mathcal O x^2\oplus\mathcal O xy\oplus\mathcal O y^2$
in the space of binary quadratic forms. The Fock Hermitian form makes
these monomials orthogonal with squared norms $2,1,2$. Dividing the
form by the squared norm $2$ of the extremal monomial $x^2$ gives
Gram matrix $\operatorname{diag}(1,1/2,1)$, and direct calculation gives
\[
 N^\#=\mathcal O x^2\oplus2\mathcal O xy\oplus\mathcal O y^2
                  \subset N.
\]
For the unnormalized Fock form, however, the dual is
$\tfrac12\mathcal O x^2\oplus\mathcal O xy
 \oplus\tfrac12\mathcal O y^2$, which is not contained in $N$.
Indeed, its pairing with $x^2$ shows that $x^2/2$ lies in the dual,
whereas $x^2/2\notin N$.
Thus even this elementary coefficient space changes the admissible
dual exponent when the normalizing scalar is removed. This is a
counterexample to that simplification, not to the normalized
construction above.
\end{example}

\subsection{Quotient metrics and the Gram determinant bound}

Fix an embedding $v_0:K\hookrightarrow\mathbb C$ with
$v_0(\bar a)=\overline{v_0(a)}$ for every $a\in K$.
Let $H=H^*\in\operatorname{GL}_n(K)$ satisfy
\begin{equation}\label{div:H-condition}
 v(\bar H)=\overline{v(H)},\qquad v(H)>0
       \quad\text{for every complex embedding }v.
\end{equation}
The first condition holds, for example, when the entries lie in a
subfield on which all embeddings intertwine the involution with complex
conjugation. Let $P\in\operatorname{Mat}_{D\times n}(K)$ have rank $D$.
Define
\begin{equation}\label{div:quotient-matrices}
 G_v=\bigl(v(P)v(H)^{-1}v(P)^*\bigr)^{-1},\qquad
 \Gamma=(PH^{-1}P^*)^{-1}.
\end{equation}
The algebraic inverse defining $\Gamma$ exists: at $v_0$ its inverse
matrix is positive definite, hence has nonzero determinant. Also
$\Gamma=\Gamma^*$.

\begin{lemma}[The quotient metric]\label{div:quotient-metric}
For a positive definite Hermitian matrix $H_{\mathbb C}$ and a surjective
matrix $P_{\mathbb C}$, put
$G=(P_{\mathbb C} H_{\mathbb C}^{-1}P_{\mathbb C}^*)^{-1}$. Then
\[
 y^*Gy=\min_{P_{\mathbb C}x=y}x^*H_{\mathbb C}x.
\]
The unique minimizing vector is
$x_0=H_{\mathbb C}^{-1}P_{\mathbb C}^*Gy$ and is orthogonal to
$\ker P_{\mathbb C}$. A linear functional vanishing on this kernel
therefore has exactly the same operator norm before and after descent
to the quotient.
\end{lemma}
\begin{proof}
We have $P_{\mathbb C}x_0=y$, and for $z\in\ker P_{\mathbb C}$,
$z^*H_{\mathbb C}x_0=z^*P_{\mathbb C}^*Gy=0$.
Every other lift is $x_0+z$, whose squared norm is
$x_0^*H_{\mathbb C}x_0+z^*H_{\mathbb C}z\ge x_0^*H_{\mathbb C}x_0=y^*Gy$,
with equality only for $z=0$. The quotient map decreases norms, while
every vector has a lift of the same norm. These two facts give the
two inequalities between the operator norms of a functional and its
descent.
\end{proof}

\begin{proposition}[A local bound from the ambient dual]\label{div:local-bound}
Let $M\subset K^n$ be a full lattice, put $\Lambda=PM$, and let
$a\in\mathcal O\setminus\{0\}$ satisfy $aM^\#\subset M$, where the
dual is for $H$. Then, with the Hermitian dual of $\Lambda$ taken for
$\Gamma$,
\[
 a\Lambda^\#\subset\Lambda,
 \qquad
 \operatorname{ord}_{\mathfrak p}
       g_\Gamma(e_{\mathfrak p})
       \le D\operatorname{ord}_{\mathfrak p}(a)
       \quad\text{for every }\mathfrak p.
\]
In particular $a=\ell^g$, $\ell\ge1$, gives the bound
$gD\operatorname{ord}_{\mathfrak p}(\ell)$.
\end{proposition}
\begin{proof}
The algebraic map $J=H^{-1}P^*\Gamma$ satisfies $PJ=1$ and
\[
 h_H(x,Jy)=h_\Gamma(Px,y).
\]
Thus $y\in\Lambda^\#$ implies $Jy\in M^\#$. The hypothesis gives
$aJy\in M$, and applying $P$ gives $ay\in\Lambda$.
At $\mathfrak p$ the inclusion $a\Lambda^\#\subset\Lambda$ yields
\[
 D\operatorname{ord}_{\mathfrak p}(a)
       +a_{\mathfrak p}(\Lambda^\#)
       \ge a_{\mathfrak p}(\Lambda).
\]
Use Lemma~\ref{div:hermitian-index} to obtain the stated bound.
\end{proof}

\begin{lemma}[Determinant Cauchy--Schwarz]\label{div:det-cs}
For $D\times n$ complex matrices $U,V$,
\[
 |\det(UV^*)|^2\le\det(UU^*)\det(VV^*).
\]
\end{lemma}
\begin{proof}
The $D$th exterior power of $\mathbb C^n$ has the Hermitian inner
product determined by
$\langle u_1\wedge\cdots\wedge u_D,
v_1\wedge\cdots\wedge v_D\rangle
=\det(\langle u_i,v_j\rangle)$.
Apply ordinary Cauchy--Schwarz to the exterior products of the rows of
$U$ and $V$. Their squared norms are the two Gram determinants on the
right. This also proves the assertion when the row rank is deficient.
\end{proof}

\begin{theorem}[Covolume of an algebraic quotient]\label{div:quotient-covolume}
Under \eqref{div:H-condition}, for a full lattice $M\subset K^n$ and
$\Lambda=PM$, choose a simultaneous local basis $e_{\mathfrak p}$ of
$\Lambda$ at each conjugate pair of primes. Then
\begin{equation}\label{div:quotient-inequality}
 c_G(\Lambda)\le
 \sum_{\mathfrak p}w_{\mathfrak p}
       \operatorname{ord}_{\mathfrak p}g_\Gamma(e_{\mathfrak p}).
\end{equation}
The sum has finite support. The difference between the right and left
sides is independent of $M$ and is nonnegative. If the involution acts
as complex conjugation at every embedding of $K$, equality holds.
\end{theorem}
\begin{proof}
On the set of all complex embeddings define
$\tau(v)(a)=\overline{v(\bar a)}$. This is an involution commuting
with conjugation of embeddings. By \eqref{div:H-condition},
$\tau(v)(H)=v(H)$. Set
\[
 A_v=v(P)v(H)^{-1}v(P)^*,\qquad
 B_v=v(P)v(H)^{-1}\tau(v)(P)^*.
\]
Thus $G_v=A_v^{-1}$ and $v(\Gamma)=B_v^{-1}$.
The matrices $B_v$ are invertible, because the algebraic matrix inverted
in \eqref{div:quotient-matrices} is invertible. Apply
Lemma~\ref{div:det-cs} to
$U=v(P)v(H)^{-1/2}$ and $V=\tau(v)(P)v(H)^{-1/2}$:
\[
 |\det B_v|^2\le\det A_v\det A_{\tau(v)}.
\]
After inversion and logarithms this becomes
\[
 \log\det G_v+\log\det G_{\tau(v)}
       \le2\log|v(\det\Gamma)|.
\]
Sum over all embeddings. Since $\tau$ is a permutation, and conjugate
embeddings give equal real terms, we obtain
\begin{equation}\label{div:quotient-arch}
 \sum_{v\mid\infty}\log\det G_v
       \le\sum_{v\mid\infty}\log|v(\det\Gamma)|.
\end{equation}
By \eqref{div:gram-valuation}, the proposed upper bound equals
\[
 \sum_{\mathfrak p}w_{\mathfrak p}
       \operatorname{ord}_{\mathfrak p}\det\Gamma
       +2\sum_{\mathfrak p}w_{\mathfrak p}a_{\mathfrak p}(\Lambda).
\]
Both sums have finite support. Apply \eqref{div:product-formula} to
$\det\Gamma$ and then \eqref{div:quotient-arch}. This is exactly
\eqref{div:quotient-inequality}, by \eqref{div:covolume}.

More explicitly the difference is
\begin{equation}\label{div:quotient-defect}
 \delta(P,H)=\frac1{2d}\sum_{v:K\hookrightarrow\mathbb C}
 \log\frac{\det A_v\det A_{\tau(v)}}{|\det B_v|^2}\ge0.
\end{equation}
This expression contains no lattice. If $\tau$ fixes every embedding,
then $A_v=B_v$, so every summand is zero.
\end{proof}

\begin{example}[Strictness over a field which is not CM]\label{div:strict}
Let $\alpha=\sqrt[3]{2}$, let $\zeta$ be a primitive cube root of unity,
and take $K=\mathbb Q(\alpha,\zeta)$ with its usual complex conjugation
at the embedding where $\alpha>0$. Put $H=I_2$, $M=\mathcal O_K^2$,
and $P=(1,\alpha)$. Then $\Lambda=\mathcal O_K$ and
\[
 c_G(\Lambda)=-\log(1+2^{2/3}),\qquad
 \sum_{\mathfrak p}w_{\mathfrak p}
       \operatorname{ord}_{\mathfrak p}g_\Gamma(1)
       =-\tfrac13\log5.
\]
The inequality is strict.
\end{example}
\begin{proof}
The field has degree six, no real embeddings, and three complex places.
At every embedding $|v(\alpha)|=\alpha$, so
$G_v=(1+\alpha^2)^{-1}$. Algebraically
$\Gamma=(1+\alpha^2)^{-1}$, but its other embeddings are not all equal
to these positive metric matrices. Since
\[
 \prod_{k=0}^2(1+\alpha^2\zeta^{2k})=1+\alpha^6=5,
 \qquad |N_{K/\mathbb Q}(1+\alpha^2)|=25,
\]
the two formulas follow from \eqref{div:covolume} and
\eqref{div:product-formula}. Strictness follows from
$(1+\alpha^2)^3>1+\alpha^6=5$.
The fixed field of conjugation is $\mathbb Q(\alpha)$, which is not
totally real; hence $K$ is not a CM field. This example satisfies all
hypotheses of Theorem~\ref{div:quotient-covolume}. It disproves replacing
its inequality by an equality in the stated generality.
\end{proof}

\subsection{Patching the lattices and an exact divisor formula}

For a $K$-linear isomorphism $f$ between $q$-dimensional spaces with
metrics $G_v,G'_v$, write
\[
 d_v(f)=\frac{|\det v(f)|^2\det G'_v}{\det G_v},\qquad
 A(f)=\frac2d\sum_{v\mid\infty}\log d_v(f).
\]
These are coordinate independent. Directly from the product formula,
\begin{equation}\label{div:image-covolume}
 c_{G'}(fL)=c_G(L)+A(f).
\end{equation}

Here are the data for the formula. They are explicit finite-dimensional
objects, rather than an assumed divisor inequality.
There is a parameter set $\mathcal T$ with a distinguished element $0$.
For each $t$ there is a $D$-dimensional Hermitian $K$-space $Q_t$, a full
lattice $\Lambda_t\subset Q_t$, and a $K$-linear isomorphism
$\iota_t:Q_t\to E$ to one fixed Hermitian space $E$.
The space $Q_t$ and lattice $\Lambda_t$ arise from matrices and lattices
as in \eqref{div:quotient-matrices}: $\Lambda_t=P_tM_t$ with quotient
metrics $G_{t,v}$ and algebraic form $\Gamma_t$.
It is enough to require the global condition
\eqref{div:H-condition} for $t=0$; for other $t$ only the equality
$v_0(\Gamma_t)=G_{t,v_0}$, nondegeneracy and Hermitian symmetry are
needed below.

For each $j$ in a finite set $J$, there is a fixed Hermitian space $W_j$
of dimension $d_j$ and a full lattice $M_j\subset W_j$. For every $t$
there is a Hermitian space $T_{j,t}$ of dimension $d_j$, with full lattice
$L_{j,t}$, algebraic nondegenerate Hermitian form $H_{j,t}$, and
isomorphism $\theta_{j,t}:T_{j,t}\to W_j$.
At $v_0$ its metric is $v_0(H_{j,t})$. For $t=0$ the metric is
$v(H_{j,0})$ at every place and the entries of $H_{j,0}$ satisfy the
conjugation condition in \eqref{div:H-condition}.

Let $\kappa_v>0$, and let $\rho_t\in K^\times$ satisfy
$\rho_0=1$ and $v_0(\rho_t)>0$. The compatibility conditions on
determinants are
\begin{align}
 d_v(\iota_0)&=\kappa_v^D\prod_{j\in J}d_v(\theta_{j,0})
                       &&(v\mid\infty),\label{div:det-base}\\
 d_{v_0}(\iota_t)&=(\kappa_{v_0}v_0(\rho_t))^D
                         \prod_{j\in J}d_{v_0}(\theta_{j,t})
                       &&(t\in\mathcal T).\label{div:det-local}
\end{align}
Below we also prove the spectral calculation which produces such
determinant identities.

Choose parameters $t_{\mathfrak p}$ such that
\begin{equation}\label{div:patch-parameters}
 t_{\overline{\mathfrak p}}=t_{\mathfrak p},\qquad
 t_{\mathfrak p}=0\text{ except at finitely many primes},\qquad
 \operatorname{ord}_{\mathfrak p}(\rho_{t_{\mathfrak p}})=0.
\end{equation}
Put $R_t=\iota_t\Lambda_t$. There is a unique full lattice $R\subset E$
with $R_{\mathfrak p}=(R_{t_{\mathfrak p}})_{\mathfrak p}$ for every
$\mathfrak p$. For completeness, choose a nonzero integer $a$ such that
$aR_0\subset R_t\subset a^{-1}R_0$ for the finitely many lattices in
question. Enlarge the exceptional set to include the primes over $a$.
The finite module $a^{-1}R_0/aR_0$ is the direct sum of its primary
components. In each component select the submodule specified by the
desired local lattice, and take its inverse image in $a^{-1}R_0$.
This constructs $R$ and proves its asserted localizations. Uniqueness
follows because a finitely generated torsion module is zero if all its
localizations are zero.

Choose simultaneous local bases $e_{t,\mathfrak p}$ of $\Lambda_t$ and
$f_{j,t,\mathfrak p}$ of $L_{j,t}$. Set
\begin{align*}
 g_{t,\mathfrak p}&=g_{\Gamma_t}(e_{t,\mathfrak p}),\\
 I_{j,t,\mathfrak p}&=i_{\mathfrak p}(M_j,\theta_{j,t}L_{j,t}),
       &J_{j,t,\mathfrak p}&=j_{\mathfrak p}(H_{j,t},L_{j,t}),\\
 F_{t,\mathfrak p}&=\operatorname{ord}_{\mathfrak p}g_{t,\mathfrak p}
       +\sum_{j\in J}\bigl(I_{j,t,\mathfrak p}
              +I_{j,t,\overline{\mathfrak p}}+J_{j,t,\mathfrak p}\bigr),
       &\overline{\log\kappa}&=\frac1{r_2}\sum_{v\mid\infty}\log\kappa_v.
\end{align*}
The letter $F$ here denotes an explicitly defined local number, not an
additional hypothesis.

\begin{theorem}[Exact divisor formula]\label{div:patch-formula}
With the preceding data,
\begin{equation}\label{div:exact-divisor}
 c_E(R)=D\,\overline{\log\kappa}+\sum_{j\in J}c_{W_j}(M_j)
     +\sum_{\mathfrak p}w_{\mathfrak p}F_{t_{\mathfrak p},\mathfrak p}
     -\delta(P_0,H_0),
\end{equation}
where $\delta(P_0,H_0)\ge0$ is given explicitly by
\eqref{div:quotient-defect}. All sums over finite primes have finite
support. Dropping this nonnegative last term gives the divisor inequality.
\end{theorem}
\begin{proof}
Fix $\mathfrak p$ and put $t=t_{\mathfrak p}$. Let $\gamma$ be the
change-of-basis matrix from $\iota_0e_{0,\mathfrak p}$ to
$\iota_te_{t,\mathfrak p}$, and let $\delta_j$ be the corresponding
matrix from $\theta_{j,0}f_{j,0,\mathfrak p}$ to
$\theta_{j,t}f_{j,t,\mathfrak p}$.
Write $h_{j,t}=g_{H_{j,t}}(f_{j,t,\mathfrak p})$.
Changing Gram determinants at $v_0$ multiplies them by
$|v_0(\det\gamma)|^2$ and $|v_0(\det\delta_j)|^2$, respectively.
Apply \eqref{div:det-base} at $v_0$ and
\eqref{div:det-local}, substitute these change-of-basis identities,
and cancel the nonzero common factors. The resulting identity is
\[
 \overline{\det\gamma}\det\gamma\,g_{0,\mathfrak p}
                  \prod_jh_{j,t}
 =\rho_t^D g_{t,\mathfrak p}
                  \prod_j\bigl(\overline{\det\delta_j}
                                    \det\delta_j\,h_{j,0}\bigr).
\]
It is an identity in $K$: its image under the injective embedding $v_0$
has just been proved, and $v_0$ respects the involution.

Take $\operatorname{ord}_{\mathfrak p}$, use the last condition of
\eqref{div:patch-parameters}, and apply Lemmas~\ref{div:index} and
\ref{div:hermitian-index}. Because the bases are simultaneous at the
two primes, this gives
\begin{equation}\label{div:local-difference}
 F_{t,\mathfrak p}-F_{0,\mathfrak p}
 =i_{\mathfrak p}(R_0,R_t)
       +i_{\overline{\mathfrak p}}(R_0,R_t).
\end{equation}
By the first condition of \eqref{div:patch-parameters}, the two local
indices on the right are also those of $R_0,R$.
The terms $F_{0,\mathfrak p}$ have finite support, by
\eqref{div:gram-valuation} and the finite support of determinant ideals.
Equation \eqref{div:local-difference} and Lemma~\ref{div:index} therefore
give finite support for the selected terms as well. Summing and pairing
conjugate primes yields
\begin{equation}\label{div:patch-sum}
 \sum_{\mathfrak p}w_{\mathfrak p}F_{t_{\mathfrak p},\mathfrak p}
 =\sum_{\mathfrak p}w_{\mathfrak p}F_{0,\mathfrak p}
       +c_E(R)-c_E(R_0).
\end{equation}

At the base parameter, \eqref{div:rational-hermitian},
\eqref{div:index-covolume} and \eqref{div:image-covolume} give
\begin{align*}
 &\sum_{\mathfrak p}w_{\mathfrak p}
     (I_{j,0,\mathfrak p}+I_{j,0,\overline{\mathfrak p}}
                                   +J_{j,0,\mathfrak p})\\
 &\hspace{10mm}=c_{W_j}(\theta_{j,0}L_{j,0})-c_{W_j}(M_j)
                                     -c_{T_{j,0}}(L_{j,0})
  =A(\theta_{j,0})-c_{W_j}(M_j).
\end{align*}
Moreover \eqref{div:det-base} implies
$A(\iota_0)=D\overline{\log\kappa}+\sum_jA(\theta_{j,0})$.
Substitute these two identities and
$c_E(R_0)=c_{Q_0}(\Lambda_0)+A(\iota_0)$ into
\eqref{div:patch-sum}. The remaining difference
$\sum_{\mathfrak p}w_{\mathfrak p}
\operatorname{ord}_{\mathfrak p}g_{0,\mathfrak p}
-c_{Q_0}(\Lambda_0)$ is exactly $\delta(P_0,H_0)$ by
Theorem~\ref{div:quotient-covolume}. This proves the formula.
\end{proof}

\begin{corollary}[Insertion of local estimates]\label{div:scalar-divisor}
Assume $D>0$, and let $S$ contain the support of the selected local
terms. Abbreviate $I_{j,\mathfrak p}=I_{j,t_{\mathfrak p},\mathfrak p}$
and $J_{j,\mathfrak p}=J_{j,t_{\mathfrak p},\mathfrak p}$.
Suppose that real numbers satisfy
\begin{align*}
 \operatorname{ord}_{\mathfrak p}g_{t_{\mathfrak p},\mathfrak p}
       &\le Dg_{\mathfrak p},
 &\sum_{\mathfrak p\in S}w_{\mathfrak p}g_{\mathfrak p}&\le L_G,\\
 \sum_j(I_{j,\mathfrak p}+I_{j,\overline{\mathfrak p}})
       &\le D\left(\sum_j(2f_{j,\mathfrak p}+e_{j,\mathfrak p})
                                      +2/\ell_{\mathfrak p}^2\right),\\
 J_{j,\mathfrak p}&\le Dq_{j,\mathfrak p}+\eta_{j,\mathfrak p},
 &c_{W_j}(M_j)/D&\le L_j,\\
 \sum_{\mathfrak p\in S}w_{\mathfrak p}
                     (e_{j,\mathfrak p}+q_{j,\mathfrak p})&\le E_{1j},
 &\sum_{\mathfrak p\in S}w_{\mathfrak p}\eta_{j,\mathfrak p}/D&\le E_{2j},\\
 \sum_{\mathfrak p\in S}2w_{\mathfrak p}/\ell_{\mathfrak p}^2&\le E_3.
\end{align*}
Then
\begin{equation}\label{div:final-divisor}
 2\mu(R^\vee)\le\overline{\log\kappa}+\sum_jL_j
   +2\sum_j\sum_{\mathfrak p\in S}w_{\mathfrak p}f_{j,\mathfrak p}
   +L_G+\sum_j(E_{1j}+E_{2j})+E_3.
\end{equation}
\end{corollary}
\begin{proof}
Divide \eqref{div:exact-divisor} by $D$, discard $-\delta/D\le0$,
and insert the displayed estimates. Every finite-prime weight is
nonnegative, so the directions of all inequalities are preserved.
Interchange the two finite sums and use \eqref{div:ordinary-dual}.
\end{proof}

\subsection{A uniform estimate for the elementary prime error}

The term $E_3$ in the preceding corollary has a uniform numerical bound
which can be proved without any prime-distribution theorem.

\begin{proposition}\label{div:prime-error}
For a prime ideal $\mathfrak p$ above the rational prime $p$, let
$e_{\mathfrak p}=\operatorname{ord}_{\mathfrak p}(p)$ and
$\ell_{\mathfrak p}=p/\sqrt{e_{\mathfrak p}}$. Then
\[
 \sum_{\mathfrak p}\frac{2w_{\mathfrak p}}{\ell_{\mathfrak p}^2}
 =2\sum_{p\ {
m prime}}\frac{\log p}{p^2}
 <\frac{99841}{100000}<1.
\]
Every sum over a finite set of prime ideals obeys the same upper bound.
\end{proposition}
\begin{proof}
Write $N\mathfrak p=p^{f_{\mathfrak p}}$. The identity
$\sum_{\mathfrak p\mid p}e_{\mathfrak p}f_{\mathfrak p}=d$ gives the
first equality; all summands are positive.
For the numerical bound let $f(x)=\log x/x^2$, which is decreasing for
$x\ge2$. Every prime larger than $200$ is $1$ or $5$ modulo $6$.
For each of these two arithmetic progressions, comparison with its
first term and the subsequent integral gives
\[
 \sum_{p>200}f(p)
 \le2f(200)+\frac13\int_{200}^{\infty}f(x)\,dx
 =\frac{2\log200}{200^2}+\frac{\log200+1}{600}
 <\frac{2153}{200000},
\]
using $\log200<53/10$.
The finite sum is bounded by the following rational table. The last
column is an upper bound in units of $10^{-6}$.
\begin{center}\small
\begin{tabular}{lll}
\toprule
Interval&Primes&$10^6\sum\log p/p^2$\quad upper bound\\
\midrule
$[2,40]$&$2,3,5,7,11,13,17,19,23,29,31,37$&$468542$\\
$[41,80]$&$41,43,47,53,59,61,67,71,73,79$&$12964$\\
$[81,120]$&$83,89,97,101,103,107,109,113$&$3757$\\
$[121,160]$&$127,131,137,139,149,151,157$&$1753$\\
$[161,200]$&$163,167,173,179,181,191,193,197,199$&$1424$\\
\bottomrule
\end{tabular}
\end{center}
Here is a purely rational way to verify every entry, including the
bound for $\log200$. For $0\le z\le1/3$ put
\[
 U(z)=2\sum_{k=0}^{15}\frac{z^{2k+1}}{2k+1}
                   +\frac{2z^{33}}{33(1-z^2)}.
\]
The series obtained by integrating
$2/(1-z^2)=2\sum_{k\ge0}z^{2k}$ term by term shows that
$\log((1+z)/(1-z))\le U(z)$: bound every denominator in the remainder
by $33$ and sum the resulting geometric series.
For an integer $n\ge1$, choose $k$ with $2^k\le n<2^{k+1}$ and set
$z=(n-2^k)/(n+2^k)$. Then
\[
 \log n\le kU(1/3)+U(z).
\]
Substitution in each entry of the table proves the displayed rational
upper bound by integer arithmetic, and substitution of $n=200$ gives
$kU(1/3)+U(z)<53/10$. The accompanying short program uses precisely
these fractions, with no floating-point comparisons.
Adding the table gives
$\sum_{p\le200}\log p/p^2<488440/10^6=12211/25000$.
Together with the tail this is $99841/200000$; multiplication by two
proves the result.
\end{proof}

\subsection{Weighted operator matrices and rectangular chains}

We next give the norm bounds used by the chain construction, including
the dependence on the number of classes. This isolates the exact
pointwise estimate that must be proved for each actual kernel.

\begin{lemma}[A weighted operator matrix]\label{div:weighted-matrix}
Let $S,T$ be finite sets, let $U_s,V_t$ be complex normed spaces, and
let $\alpha_s,\beta_t>0$. For continuous linear maps
$A_{ts}:U_s\to V_t$ define
\[
 (\mathcal A x)_t=\sum_{s\in S}\alpha_s A_{ts}x_s,
 \qquad \|x\|_\alpha^2=\sum_s\alpha_s\|x_s\|^2,
 \qquad \|y\|_\beta^2=\sum_t\beta_t\|y_t\|^2.
\]
If $\|A_{ts}\|^2\le a_{ts}^2$, with $a_{ts}\ge0$, and
$M_{ts}=\sqrt{\beta_t\alpha_s}\,a_{ts}$, then
\begin{equation}\label{div:weighted-sharp}
 \|\mathcal A\|\le\|M\|_{\ell^2(S)\to\ell^2(T)}
       \le\left(\sum_{t,s}\beta_t\alpha_s a_{ts}^2\right)^{1/2}.
\end{equation}
In particular, if $\|A_{ts}\|^2\le A$, then
\begin{equation}\label{div:weighted-coarse}
 \|\mathcal A\|^2\le A\left(\sum_s\alpha_s\right)
                           \left(\sum_t\beta_t\right).
\end{equation}
If $S=\coprod_{a\in\mathcal A_0}S_a$ and
$T=\coprod_{b\in\mathcal A_1}T_b$, with
$\alpha_s,\beta_t\le1$, $|S_a|\le h$, $|T_b|\le h'$, this yields
$\|\mathcal A\|^2\le A|\mathcal A_0||\mathcal A_1|hh'$.
\end{lemma}
\begin{proof}
Put $u_s=\sqrt{\alpha_s}\|x_s\|$. The triangle inequality gives
\[
 \sqrt{\beta_t}\| (\mathcal A x)_t\|
       \le\sum_sM_{ts}u_s.
\]
Take the Euclidean norm in $t$ and use the definition of $\|M\|$.
The operator norm of a scalar matrix is at most its Hilbert--Schmidt
norm, proving \eqref{div:weighted-sharp}. The remaining claims follow
by summing the weights and the cardinality bounds. The argument includes
empty index sets: the corresponding map and sums are zero. Factors of
absolute value one in $A_{ts}$ do not affect any bound.
\end{proof}

The first bound in \eqref{div:weighted-sharp} is sometimes substantially
better than the class-count bound. Thus no enlargement of the number of
classes is intrinsically required in this part of the argument.

Suppose now that $V_0,\ldots,V_m$ are Hermitian $K$-spaces of possibly
different dimensions, and $K_j:V_{j-1}\to V_j$ are $K$-linear maps.
At each complex place suppose there are norm-preserving realizations
$\Phi_{j,v}$ into weighted spaces as above and exact identities
\[
 \Phi_{j,v}(K_{j,v}x)=\mathcal A_{j,v}\Phi_{j-1,v}(x).
\]
Linearity or surjectivity of $\Phi_{j,v}$ is not needed for the following
conclusion; preservation of norms and the displayed identity suffice.
Lemma~\ref{div:weighted-matrix} gives $\|K_{j,v}\|^2\le C_{j,v}$ with
\[
 C_{j,v}=A_{j,v}\left(\sum_s\alpha_{j-1,s}\right)
                          \left(\sum_t\alpha_{j,t}\right),
\]
or with the sharper scalar-matrix bound. Consequently
\begin{equation}\label{div:chain}
 \|K_{m,v}\cdots K_{1,v}\|^2\le\prod_{j=1}^m C_{j,v}.
\end{equation}
This follows by applying the bound successively to a vector. It requires
no square matrices, determinant nonvanishing, or equality of dimensions.
The empty chain is the identity.

\begin{lemma}[Spectral determinant and lower norm]\label{div:spectral}
Let $\iota:Q\to E$ be an isomorphism of finite-dimensional complex
Hermitian spaces, and let a finite nonempty family of subspaces $Q_p$
span $Q$. Suppose
$s_p>0$ and
\begin{equation}\label{div:spectral-cross}
 \langle\iota x,\iota y\rangle_E
       =s_p\langle x,y\rangle_Q
             \quad(x\in Q_p,\ y\in Q).
\end{equation}
Then
$\|\iota x\|_E^2\ge(\min_ps_p)\|x\|_Q^2$ for every $x\in Q$.
If the sum of the $Q_p$ is direct and $d_p=\dim Q_p$, then
\[
 d(\iota)=\prod_ps_p^{d_p}.
\]
In particular, if $s_p=\kappa\prod_j\lambda_{j,p}$ and
$\theta_j:T_j\to W_j$ has a direct-sum decomposition into subspaces
of dimensions $d_p$ on which its pullback metric has scalar
$\lambda_{j,p}$ in the sense of \eqref{div:spectral-cross}, then
\[
 d(\iota)=\kappa^{\dim Q}\prod_jd(\theta_j).
\]
\end{lemma}
\begin{proof}
Let $B=\iota^*\iota$, the positive self-adjoint operator on $Q$.
Equation \eqref{div:spectral-cross}, for all $y$, says exactly that
$Bx=s_px$ on $Q_p$. Since these spaces span $Q$, all eigenvalues of
$B$ occur among the $s_p$, proving the lower norm estimate by the
spectral theorem. If the sum is direct, concatenate bases of its
summands; the matrix of $B$ is diagonal by blocks, so
$\det B=\prod_ps_p^{d_p}$. This is the determinant ratio. The final
claim follows by substituting $s_p$ and applying the same calculation
to each $\theta_j$.
\end{proof}

\begin{proposition}[Height of a descended pairing]\label{div:pairing-height}
Let $P:V\to Q$ be surjective and give $Q$ its quotient metric. Suppose
$\iota:Q\to E$ satisfies Lemma~\ref{div:spectral} with
$s_{p,v}=\kappa_v\prod_j\lambda_{j,p,v}>0$ at each infinite place.
Let $K_m\cdots K_1:U\to U_m$ be a rectangular chain and let
$\mathcal P:U_m\times V\to K$ and $\mathcal B:U\times E\to K$
be bilinear pairings satisfying the exact identity
\begin{equation}\label{div:pairing-factorization}
 \mathcal B(x,\iota(Pz))
       =\mathcal P(K_m\cdots K_1x,z)
                         \quad(x\in U,z\in V).
\end{equation}
If $\|\mathcal P_v\|^2\le e^{h_{P,v}}$ and
$\|K_{j,v}\|^2\le e^{h_{j,v}}$, then
\begin{equation}\label{div:pairing-bound}
 \|\mathcal B_v\|^2
 \le\exp\left(h_{P,v}+\sum_jh_{j,v}
            -\log\kappa_v-\min_p\sum_j\log\lambda_{j,p,v}\right).
\end{equation}
For $\mathcal B\ne0$, define its squared-norm height by
$h_2(\mathcal B)=r_2^{-1}\sum_v\log\|\mathcal B_v\|^2$.
Then averaging \eqref{div:pairing-bound} gives
\[
 h_2(\mathcal B)\le\frac1{r_2}\sum_v
       \left(h_{P,v}+\sum_jh_{j,v}
                   -\min_p\sum_j\log\lambda_{j,p,v}\right)
       -\overline{\log\kappa}.
\]
\end{proposition}
\begin{proof}
Fix $v$, $x\in U_v$, $e\in E_v$, and put $q=\iota_v^{-1}e$.
By Lemma~\ref{div:quotient-metric} there is a lift $z$ of $q$ with
$\|z\|_V=\|q\|_Q$. Lemma~\ref{div:spectral} gives
\[
 \|q\|_Q\le\frac{\|e\|_E}
           {\sqrt{\kappa_v\min_p\prod_j\lambda_{j,p,v}}}.
\]
Extend \eqref{div:pairing-factorization} by complex bilinearity and
use the operator bounds and submultiplicativity:
\[
 |\mathcal B_v(x,e)|
 \le\frac{\|\mathcal P_v\|\prod_j\|K_{j,v}\|}
           {\sqrt{\kappa_v\min_p\prod_j\lambda_{j,p,v}}}
          \|x\|_U\|e\|_E.
\]
Square and use positivity to take logarithms of the denominator.
Since logarithm is increasing, the logarithm of the minimum of these
positive products is the minimum of their sums of logarithms.
If $\mathcal B\ne0$, a nonzero matrix entry remains nonzero under every
field embedding, so all its local operator norms are positive and the
height is defined. Averaging proves the last assertion. If
$\mathcal B=0$, the norm bound still holds, but no logarithm of zero is
needed.
\end{proof}

To use these bounds in the ordinary-slope comparison of
Section~\ref{cs-section}, the conventions must be converted together.
The bilinear pairing $\mathcal B$ and the induced linear map
$U\to E^\vee$ have the same operator norm. Since $n_v=2$ at a complex
place and $r_2=d/2$, the height used there satisfies
\[
 h_2(\mathcal B)=2h(\mathcal B),\qquad
 c(R)/D=2\mu(R^\vee).
\]
Accordingly, add the divisor inequality
\eqref{div:final-divisor} to the squared-norm height inequality above,
cancel the two occurrences of $\overline{\log\kappa}$, and divide
the whole inequality by two before applying the ordinary-slope
comparison. Converting only one of the two inequalities would give
an incorrect factor of two.

\subsection{Local divisibility gives an integral map}

\begin{proposition}\label{div:integral-pairing}
Let $N\subset U$ and $R\subset E$ be full lattices, with $\dim U>0$, and let
$\mathcal B:U\times E\to K$ be bilinear. For each prime let
$s_{\mathfrak p}\in\mathbb R$, with
$k_{\mathfrak p}=\lfloor s_{\mathfrak p}\rfloor$ of finite support.
Assume that for all $x\in N$, all $y\in R_{\mathfrak p}$, and every
prime $\mathfrak p$, either $\mathcal B(x,y)=0$ or
\[
 \operatorname{ord}_{\mathfrak p}\mathcal B(x,y)\ge s_{\mathfrak p}.
\]
Put $\mathfrak b=\prod_{\mathfrak p}\mathfrak p^{-k_{\mathfrak p}}$.
Then $\mathcal B$ induces an integral linear map
$\mathfrak bN\to R^\vee$, and
\[
 2\mu(\mathfrak bN)=2\mu(N)
             +2\sum_{\mathfrak p}w_{\mathfrak p}k_{\mathfrak p}.
\]
It is enough to establish the valuation hypothesis on any lattice with
the same localization as $R$ at the prime in question.
\end{proposition}
\begin{proof}
At $\mathfrak p$, denominators used to localize $N$ are units. Thus the
hypothesis extends to $x\in N_{\mathfrak p}$ by linearity and the
valuation inequality for sums. Multiplication by
$\mathfrak b_{\mathfrak p}=\pi^{-k_{\mathfrak p}}\mathcal O_{\mathfrak p}$
lowers a valuation by $k_{\mathfrak p}$, leaving it at least
$\lceil s_{\mathfrak p}\rceil-\lfloor s_{\mathfrak p}\rfloor\ge0$.
Hence $\mathcal B(\mathfrak bN,R)$ lies in every
$\mathcal O_{\mathfrak p}$, and therefore in $\mathcal O$.
This is precisely the asserted integral map.
If $a=\dim U>0$, then
$\mathfrak a(\mathfrak bN)=\mathfrak b^a\mathfrak a(N)$.
Taking logarithms in \eqref{div:covolume} and dividing by $a$ proves
the slope identity. In rank zero the map assertion remains true and
there is no slope assertion. Finally the valuation condition concerns
only the localization, so replacing a lattice by one with the same
localization does not change it.
\end{proof}

\subsection{The exact denominator of a chain}
\label{shared-level-chain}

The auxiliary-prime calculation can keep each step's cost and the
positive factor already present in the source lattice.  The following
identity does not require the intermediate operators to be invertible.

\begin{proposition}\label{shared-level-chain-integrality}
Let $R$ be a discrete valuation ring with fraction field $K$, normalized
valuation $v$, and a field extension $E/K$.
Fix $q\in R\setminus\{0\}$, with $e=v(q)>0$.
For $0\le i\le m$, let $\Lambda_i$ be a finitely generated
$R$-submodule of a finite-dimensional $E$-space $V_i$, and let
$U_i\in\operatorname{End}_E(V_i)$.
Let $F_i:V_{i-1}\to V_i$ satisfy
\[
 U_iF_i=q^{r_i}F_iU_{i-1},\qquad
 A_i:=q^{d_i}U_iF_i,\qquad A_i\Lambda_{i-1}\subseteq\Lambda_i
 \quad(1\le i\le m),
\]
where $r_i,d_i\in\mathbb Z$.  If
$M\subseteq q^aU_0^m\Lambda_0$, with $a\in\mathbb Z$, put
\[
 F=F_m\cdots F_1,\qquad
 \ell=\sum_{i=1}^m d_i+
           \sum_{i=1}^m(m-i+1)r_i-a.
\]
Then $q^\ell F(M)\subseteq\Lambda_m$.  More generally,
$cF(M)\subseteq\Lambda_m$ whenever $c\in K^\times$ and
$v(c)\ge e\ell$.  Negative values of $\ell$ are allowed.
\end{proposition}
\begin{proof}
Repeatedly move an operator $U_i$ across the maps to its right.  Each
commutation with $F_i$ occurs $m-i+1$ times, giving
\[
 A_m\cdots A_1
 =q^{\sum_i d_i+\sum_i(m-i+1)r_i}F U_0^m.
\]
The left side takes $\Lambda_0$ into $\Lambda_m$.  For
$x=q^aU_0^my\in M$, this identity gives
$q^\ell F(x)=A_m\cdots A_1(y)\in\Lambda_m$.
Finally $v(cq^{-\ell})\ge0$, so multiplication by this element
preserves $\Lambda_m$.
\end{proof}

This exponent is optimal from the displayed assumptions.  In rank one,
take $\Lambda_i=R$, $U_0=1$,
$U_i=q^{r_1+\cdots+r_i}$, and $F_i=q^{-d_i}U_i^{-1}$.
Then $A_i=1$ and, for $M=q^aR$, one has
$F(M)=q^{-\ell}R$.  The empty chain also gives the asserted formula.

This formulation permits complex ambient function spaces with their
embedded integral lattices. Membership in the original rational image
supplies the vector $y$ used in the proof; each integral $A_i$ then
supplies the next rational lift. No separate rationality of $F_i$ or
$U_i$ is assumed.

For the arithmetic chain, let $T$ be the auxiliary rational prime,
$t$ odd, $h=(t+1)/2$, and $J$ the selected set of coefficient
embeddings, of size $j$.  At an active place above $T$, the prescribed
source lattice is contained in
\[
 T^a U_0^m\Lambda_0,\qquad
 a=m\sum_{\tau\in J}b_\tau.
\]
The original local kernel and value-lattice identities give the
single-step integers
\[
 d_i=fN_{i-1}+\sum_\tau s_{i-1,\tau}
                    +\sum_{\tau\notin J}\lambda_{i,\tau},
 \qquad r_i=hj.
\]
Here $s_{i-1,\tau}$ is the determinant shift to be cleared and
$\lambda_{i,\tau}$ is the first residual highest weight after that
shift.  These identities concern the actual coefficient lattices and
the actual kernels; the stabilizer orders in their finite sums must
be units at the place.  The resulting denominator is therefore
\begin{equation}\label{shared-level-chain-exact-loss}
 \ell=\sum_{i=1}^m d_i+hj\frac{m(m+1)}2
                         -m\sum_{\tau\in J}b_\tau.
\end{equation}
If the original weight bounds give $d_i\le2W+fN_{i-1}$ and
$N_i=k+i(t+1)$, then the former common-step upper bound
$B=m(2W+fN_m)+hj m(m+1)/2$ improves to
\[
 \ell\le B-\frac{f(t+1)m(m+1)}2-m\sum_{\tau\in J}b_\tau.
\]
The gain uses the same source lattice; its degree is not changed or
counted a second time.  At inactive places, the integral-step argument
still gives exponent zero.  The local kernel identities, rational
realizations, twisted-chain identifications and all other period
denominators remain required.  This computation supplies the passage
from those local identities to the full-chain bound, not their
arithmetic construction.


\subsection{Tests of the hypotheses and the remaining arithmetic work}

The following examples check three hypotheses which cannot be omitted.
They are counterexamples to weakened statements, not counterexamples
to the propositions proved above.

\begin{example}[One prime does not suffice for a Hermitian basis]
Take $K=\mathbb Q(i)$, $L=\mathbb Z[i]$, $H=(1)$,
$\mathfrak p=(2+i)$ and $\overline{\mathfrak p}=(2-i)$.
The vector $e=2-i$ is a basis of $L_{\mathfrak p}$, because it is a unit
there. Its Gram determinant is $5$, of valuation one at $\mathfrak p$.
But $L^\#=L$, so $j_{\mathfrak p}(H,L)=0$.
Thus \eqref{div:j-gram} would be false with an ordinary local basis at
only $\mathfrak p$. The simultaneous-basis hypothesis excludes this
example exactly because $e$ is not a unit at
$\overline{\mathfrak p}$.
\end{example}

\begin{example}[Bounds on spanning lines do not suffice]
Let $Q=E=\mathbb C^2$ with standard metrics,
$\iota=\operatorname{diag}(1,\eta)$, where $0<\eta<1$, and let
$Q_1=\mathbb Ce_1$, $Q_2=\mathbb C(e_1+\varepsilon e_2)$ with
$\varepsilon>0$. These lines span $Q$, and the squared norm ratios
on them are $1$ and $(1+\eta^2\varepsilon^2)/(1+\varepsilon^2)$.
Both are at least $(1+\varepsilon^2)^{-1}$, whereas the smallest
squared singular value of $\iota$ is $\eta^2$, which can be arbitrarily
small. The all-$y$ condition in \eqref{div:spectral-cross} is therefore
essential; knowing the norm identity only within each spanning
subspace would not prove Lemma~\ref{div:spectral}.
\end{example}

\begin{example}[Indices cannot be bounded above by multiplying minors]
Over a discrete valuation ring with uniformizer $\pi$, take
$A=(1,1)$ and $B=(1,-1+\pi^N)^t$. The first determinantal ideals of
$A$ and $B$ are both the unit ideal, while $AB=(\pi^N)$ has index
$N$. This holds for every $N\ge1$.
Thus Lemma~\ref{div:minors} gives a lower bound on the valuation of
the product ideal, and supplies no upper bound for it. The construction
must control the actual product or the actual local pairings when an
upper bound is needed.
\end{example}

The divisor inequality and the chain estimate have now been proved from
ordinary linear algebra, local algebra, the product formula and one
explicit finite rational calculation. The non-CM quotient issue,
the conjugate primes, all normalizing factors, and rectangular matrix
dimensions are accounted for. The construction required for the GRH
argument still has to supply the following facts for its particular
family, uniformly in the large parameters:
\begin{enumerate}
\item the ambient forms $H_t$, lattices $M_t$, quotient maps $P_t$,
      comparison maps $\iota_t,\theta_{j,t}$ and the simultaneous
      spectral decompositions proving
      \eqref{div:det-base}--\eqref{div:det-local};
\item the underlying unimodular arithmetic lattices that enter the
      coefficient construction, together with the remaining local
      index estimates and their error sums in
      Corollary~\ref{div:scalar-divisor};
\item the exact matrix factorization
      \eqref{div:pairing-factorization}, the isometric realizations of
      its individual steps, and bounds for the actual operator entries
      $A_{ts}$ and the final pairing $\mathcal P$;
\item the valuation inequality of
      Proposition~\ref{div:integral-pairing} for all vectors of the
      actual source and target lattices, together with the lower
      spectral bounds occurring in \eqref{div:pairing-bound}.
\end{enumerate}
The coefficient duality giving the initial containment
\eqref{div:class-lattices}, with $s=1$ for unimodular input lattices,
is proved in Section~\ref{coef:section}. The estimates concerning zeros in the supplied quasi-Riemann papers
do not by themselves imply any of these lattice containments or exact
factorizations. They may improve the auxiliary-prime choices used to
construct the family. They do not remove the need to construct the
operators and prove their uniform bounds. Establishing these remaining
facts is a mathematical reconstruction task; their absence from this
section is not evidence that the formal development assumes them
without proof.


% Concrete arithmetic coefficient maps, their divisibility and their spectra.
% All labels in this fragment have the prefix am-.
